\section{Results and Research Decisions}
\subsection{Verified historical starting point}
Table~\ref{tab:historical} reports independently accepted completions, not raw position hits. Ordinary residual PPO improves unseen-geometry aggregate completion from feedback's 25\% to a seed mean of 56.25\%, with sample standard deviation 15.63 percentage points. This is preliminary residual adaptation evidence, not a novel-method result. ID and dynamics results are close to feedback.

\begin{table*}[t]\centering\small
\caption{Historical v6 independently accepted success generated from audited case records. Each split has 32 matched cases, comprising 16 P1 and 16 P2 cases; triples are seeds 31/41/51, not independent case replications. P1 cannot substitute for P2 evidence. These are not new pilot results.}\label{tab:historical}
\begin{tabular}{llrrr}
\toprule
Method & Task & ID & Geometry & Dynamics \\
\midrule
Damped feedback & All & 29 / 32 & 8 / 32 & 28 / 32 \\
 & P1 & 16 / 16 & 3 / 16 & 16 / 16 \\
 & P2 & 13 / 16 & 5 / 16 & 12 / 16 \\
\addlinespace
Direct PPO & All & 0, 0, 0 / 32 & 0, 0, 0 / 32 & 0, 0, 0 / 32 \\
 & P1 & 0, 0, 0 / 16 & 0, 0, 0 / 16 & 0, 0, 0 / 16 \\
 & P2 & 0, 0, 0 / 16 & 0, 0, 0 / 16 & 0, 0, 0 / 16 \\
\addlinespace
Residual PPO & All & 28, 28, 29 / 32 & 13, 18, 23 / 32 & 28, 29, 29 / 32 \\
 & P1 & 16, 16, 16 / 16 & 7, 12, 16 / 16 & 16, 16, 16 / 16 \\
 & P2 & 12, 12, 13 / 16 & 6, 6, 7 / 16 & 12, 13, 13 / 16 \\
\addlinespace
Residual + hold/proceed & All & 28, 30, 30 / 32 & 15, 15, 18 / 32 & 28, 29, 29 / 32 \\
 & P1 & 16, 16, 16 / 16 & 9, 8, 12 / 16 & 16, 16, 16 / 16 \\
 & P2 & 12, 14, 14 / 16 & 6, 7, 6 / 16 & 12, 13, 13 / 16 \\
\addlinespace
\bottomrule
\end{tabular}

\end{table*}

The more informative P2 breakdown is much less favorable: residual PPO obtains 6/6/7 successes out of 16, compared with feedback's 5. P1 accounts for a large share of aggregate improvement. Across common geometry successes, residual completion is 1.10--1.31 s faster depending on seed, but this is conditioned on joint success, not an unconditional task-time advantage. All models cover both tasks; a shared-training benefit has not been measured against single-task training.

\safewidefigure{Fig4_historical_online_learning.pdf}{Actual historical online training curves. Seed traces and interaction counts are shown; training reward and frozen deterministic task success are distinct measures.}{fig:learning}
\safewidefigure{Fig5_historical_frozen_tests.pdf}{Historical frozen physical evaluations with per-seed and P1/P2 decomposition. The original tests are now descriptive evidence, not a new method-development validation set.}{fig:tests}

\subsection{Negative mechanism evidence}
Paired complete physical replays compare waiting 0, 12 and 30 control steps in a legitimate handoff region before an identical suffix controller. Among 46 common valid handoffs, suffix successes are 29, 27 and 27 while average handoff speed decreases from 28.7 to 21.2 and 7.1 mm/s. Waiting alters several state variables, so velocity is not isolated causally; the experiment does not support a general waiting benefit.

Learned hold/proceed is not consistently better than ordinary residual PPO. Forcing its seed31 policy to resume preserves both the 54/64 validation count and successful-case set. A separate active-likelihood hybrid pilot changes the actual actor surrogate, but yields only 56 versus 57/64; it fails the predeclared four-case improvement threshold. These outcomes are retained as negative evidence rather than described as successful innovations.

\subsection{Preserved v7 development pilot}
\begin{table*}[t]\centering\small
\caption{Preserved v7 P2 development pilot from actual run records, not the subsequent R1--R3 task versions. Reported successes are independently accepted. Interactions denotes the underlying policy's training budget: inference interventions do not incur that training budget again. Nonlearning probes have zero training, not zero physical evaluation cost.}\label{tab:pilot}
\begin{tabular}{llllr}
\toprule
Method & Seed & Task & Independent success & Model train steps \\
\midrule
B0 & 31 & P2 & 30/32 & 0 \\
Feedback + backoff & 31 & P2 & 30/32 & 0 \\
B1 & 31 & P2 & 30/32 & 16384 \\
B1 & 41 & P2 & 29/32 & 16384 \\
B1 & 51 & P2 & 30/32 & 16384 \\
B2 & 31 & P2 & 30/32 & 16384 \\
B2 & 41 & P2 & 30/32 & 16384 \\
B2 & 51 & P2 & 30/32 & 16384 \\
B3 measured-history mask (eval) & 31 & P2 & 30/32 & 16384 \\
B3 next-goal mask (eval) & 31 & P2 & 30/32 & 16384 \\
B3 no next goal (train) & 31 & P2 & 30/32 & 16384 \\
B3 history mask (eval) & 31 & P2 & 30/32 & 16384 \\
B3 & 31 & P2 & 30/32 & 16384 \\
B3 & 41 & P2 & 30/32 & 16384 \\
B3 & 51 & P2 & 30/32 & 16384 \\
Fixed contact +4 mm & 31 & P2 & 29/32 & 0 \\
Fixed contact -4 mm & 31 & P2 & 30/32 & 0 \\
\bottomrule
\end{tabular}

\end{table*}
\begin{table*}[t]\centering\small
\caption{Mechanism interventions and resource accounting generated from the actual pilot evidence. An input or action effect is not by itself evidence of improved completion.}\label{tab:ablation}
\resizebox{\textwidth}{!}{\begin{tabular}{lllllr}
\toprule
Variant & Frame & Result/command history & Next goal & Seeds & Train budget \\
\midrule
B0 & Task & Zero / none & Observed & 31 & 0 \\
Feedback + backoff & Task & Zero / none & Observed & 31 & 0 \\
B1 & Body & Zero / none & Observed & 31,41,51 & 16384 \\
B2 & Task & Zero / none & Observed & 31,41,51 & 16384 \\
B3 & Task & Observed & Observed & 31,41,51 & 16384 \\
B3 no measured effect (eval) & Task & Command memory only & Observed & 31 & Same checkpoint \\
B3 no next goal (train) & Task & Observed & Hidden during training & 31 & 16384 \\
B3 no next goal (eval) & Task & Observed & Suppressed at eval & 31 & Same checkpoint \\
B3 no history (eval) & Task & Suppressed at eval & Observed & 31 & Same checkpoint \\
Fixed contact +4 mm & Task & Zero / none & Observed & 31 & 0 \\
Fixed contact -4 mm & Task & Zero / none & Observed & 31 & 0 \\
\bottomrule
\end{tabular}
}
\end{table*}
The new cohort is a development split and cannot be pooled with historical tests. Each B1/B2/B3 model receives 16,384 online interactions. Independent P2 success on 32 matched validation cases is 30/29/30 for B1, 30/30/30 for B2 and 30/30/30 for B3; corrected B0 feedback achieves 30/32. The paired B3--B2 count differences are 0/0/0. Thus the preregistered requirement of a positive difference in every seed and a mean gain of at least two cases is not met. No five-seed expansion or final-test model selection is performed. The much higher new P2 completion cannot be attributed to a novel algorithm: this cohort uses different geometry groups and a repaired shared nominal controller.

For seed31, removing all eight history channels, removing measured response while retaining command memory, or partially masking next-contract inputs preserves the same 30/32 successful-case set. Removing history changes the mean per-case applied-reference-increment RMS by about 20.6 micrometres and object-trajectory RMS by 54.2 micrometres. All compared initial robot/object/tool states match exactly in the recorded reset variables. These are real action/trajectory effects, not task-completion benefits; subsequent states are not forced to match and PhysX contact caches are not claimed to be restored. A separately trained no-next seed31 model also achieves 30/32. Fixed contact-target probes do not establish a recovery advantage and cost three full executions per case rather than one.

The new budget is 147,456 interactions for nine main models plus 16,384 for a separately trained ablation (163,840 formal), and 32,768 for two retained pre-fix software diagnostics: 196,608 new online interactions and 12 trained models in total. They are not the historical 901,120 interactions. There are 624 completed formal evaluation episodes, including training-geometry checks. Actor optimizer records show actual nonzero updates; physical trajectory and task checks are separate evidence. Persistent failures occur during phase1, approximately 29--38 mm from the goal with simulated net-contact proxies of 45--59 N, rather than only failing the final stability window.

The $-4/+4$ mm contact-target probes change mean handoff object position by approximately 1.3--1.5 mm relative to zero offset; they finish 30/29 cases versus 30 at zero. The positive offset loses one case and neither probe rescues a failure. Thus a nonzero command has some physical authority but is not a measured 4 mm displacement or a useful outcome mechanism here. A bounded existing withdrawal rule triggers three attempts on each of the two persistent failures and rescues neither; this does not prove geometric impossibility or absence of other recoveries.

Common-success mean times for B2 versus B3 are 9.3022 versus 9.1333 s, 9.4650 versus 9.2161 s, and 9.3711 versus 9.3217 s for the three seeds. The nominal controller takes 9.2461 s and the same B3 checkpoint without history takes 9.1294 s. The small B3--B2 timing difference therefore does not establish an effect-history mechanism or consistent superiority to strong feedback, and it cannot replace the preregistered primary-success criterion.

\safewidefigure{Fig6_new_pilot.pdf}{Preserved v7 P2 development evaluation and same-checkpoint conditioning interventions. B3 does not improve over B2. Mask interventions are evaluations of one trained model, not additional training seeds. This is not an R1--R3 result.}{fig:newpilot}
\safewidefigure{FigS2_new_pilot_online_learning.pdf}{Preserved v7 online episode curves under its fixed pilot budget. Exploratory episode reward/success is distinct from independent validation. Later R1--R3 interaction budgets are not included here.}{fig:newlearning}

\subsection{Training correctness}
The old episode-only credit prototype modified just 23.39\% of seed31 samples; its boundary and all-stage modes coincided because eligible episodes ended in stage0. Corrected current-rollout bootstrapping has complete coverage but is not a stage-specific innovation. Exact post-step logging removes a terminal velocity timing ambiguity without relaxing thresholds. Final deterministic versus stochastic replay of an old model yields 51 versus 28/64, illustrating why early exploratory training episodes cannot be compared directly with a final mean policy. Direct PPO remains unreliable even after removing an unused action dimension; its zero final result is not evidence that all reasonable direct RL methods fail.
